\subsubsection{Partícula en trayectoria circular sobre plano inclinado}

Una partícula $A$ de masa $m$ se une a un extremo de una cuerda de longitud $R$. El otro extremo de la cuerda se fija en un plano inclinado que forma un ángulo $\phi$ con la horizontal, como se muestra en la \textbf{Figura \ref{fig:particula_cuna}}. Considerando que hay gravedad, que existe un coeficiente de fricción deslizante $\mu_e$ entre la partícula $A$ y la superficie inclinada, y que $A$ tiene rapidez $v_0$ tangencial a la circunferencia cuando está en el punto $B$, realice lo siguiente:

\begin{enumerate}
    \item Defina los marcos de referencia y sistemas de coordenadas adecuados.
    \item Analice la cinemática de la partícula en términos del ángulo $\theta$ medido sobre el plano.
    \item Analice la cinética de la partícula mediante la Segunda Ley de Newton para obtener:
    \begin{itemize}
        \item La ecuación diferencial que modela el comportamiento de $\theta$.
        \item La velocidad $v_0$ necesaria para dar exactamente 1 vuelta.
    \end{itemize}
\end{enumerate}

\begin{figure}[htbp]
    \centering
    % --- VISTA A: PERSPECTIVA DE LA CUÑA ---
    \begin{minipage}{0.48\textwidth}
        \centering
        \begin{tikzpicture}[>=Latex, scale=1.2]
            % Cuña (Triángulo base)
            \draw[thick] (0,0) -- (3.5,0) -- (3.5, 1.5) -- cycle;
            \draw (1,0) arc (0:23:1);
            \node at (0.8, 0.17) {$\phi$};
            
            % Superficie inclinada (Trayectoria elíptica por perspectiva)
         
                \draw[very thick] (0.7,0.3)--(1.75,0.75);
                \shadedraw[ball color=gray!30] (0.7,0.3) circle (2.5pt) node[above left] {$A$};
                \filldraw[black] (1.75,0.75) circle (1pt);
            % Vector Gravedad
            \draw[->, thick] (-0.3, 1.2) -- (-0.3, 0.4) node[midway, left] {$\mathbf{g}$};
        \end{tikzpicture}
        \vfill
        \small (a) Vista lateral de la cuña.
    \end{minipage}
    \hfill
    % --- VISTA B: VISTA NORMAL AL PLANO (DCL) ---
    \begin{minipage}{0.48\textwidth}
        \centering
        \begin{tikzpicture}[>=Latex, scale=1.2]
            \def\R{1.5}
            \def\angT{160}
            
            % Trayectoria circular real
            \draw[dashed, thick] (0,0) circle (\R);
            \draw[gray!50] (-\R-0.3, 0) -- (\R+0.3, 0); % Línea horizontal del plano
            
            % Centro y Partícula
            \coordinate (O) at (0,0);
            \coordinate (A) at (\angT:\R);
            \filldraw (O) circle (1.5pt);
            \draw[thick] (O) -- (A) node[midway, below left] {$R$};
            \shadedraw[ball color=gray!30] (A) circle (3.5pt) node[above left] {$m$};
            
            % Punto B (Cénit)
            \filldraw (0, \R) circle (1.5pt) node[above] {$B$};
            \filldraw (0, -\R) circle (1.5pt) node[above] {$A$};
            
            % Ángulo theta
            \draw[->] (0.6,0) arc (0:\angT:0.6);
            \node at (0.1, 0.3) {$\theta$};
            
        \end{tikzpicture}
        \vfill
        \small (b) Vista normal al plano inclinado.
    \end{minipage}

    \caption{Esquema del movimiento de la partícula $A$ sobre la cuña inclinada.}
    \label{fig:particula_cuna}
\end{figure}